\section{RT-2、RT-X、OpenVLA 与 Action Head}

前面已经说明，机器人基座模型需要真实机器人数据和跨任务策略；本章进一步进入最接近当前 VLA 主线的论文群。这里的问题变成：既然 VLM 已经从互联网图文中学到大量视觉语言知识，机器人能否把这部分知识迁移到动作输出，而不是从少量机器人数据重新学习世界常识？

本章进入当前 VLA 主线。RT-2 被视为 VLA 的开山工作之一，它把 web-scale 的视觉语言知识迁移到机器人控制中。关键思想是：一个预训练好的 VLM 已经掌握大量视觉和语言知识，能否通过机器人动作数据，让它直接输出 action？这把互联网知识和机器人控制连接起来。

\begin{figure}[H]
\centering
\includegraphics[width=0.96\textwidth]{ch23-05001.jpg}
\caption{RT-2：Vision-Language-Action Models transfer web knowledge to robotic control。投屏帧，约 01:23:21。}
\end{figure}

\begin{knowledgebox}{读图：RT-2 用 VLM backbone 输出机器人动作}
图中可以看到 RT-2 把 VLM 和机器人 action 连接起来。核心是迁移：模型不是只在机器人数据上从零学，而是把互联网视觉语言知识迁移到机器人控制。
\end{knowledgebox}

\subsection{RT-X 与 Open X-Embodiment}

本节把焦点从单个模型结构移到数据基础设施。RT-2 证明了 VLM 迁移到机器人控制有价值，但如果训练数据只来自少数机器人和少数任务，模型仍然容易被本体、场景和动作空间锁死。因此 RT-X/Open X-Embodiment 的核心问题，是怎样把分散在不同实验室、不同机器人上的经验汇成可训练的通用数据。

RT-X 和 Open X-Embodiment 强调数据集和跨本体模型。机器人数据稀缺且碎片化，单一实验室很难覆盖足够多任务和机器人本体。Open X-Embodiment 通过汇集多机器人、多任务数据，尝试训练跨本体模型，显示通用策略优于单任务专用策略的可能性。

\begin{figure}[H]
\centering
\includegraphics[width=0.96\textwidth]{ch24-05336.jpg}
\caption{RT-X / Open X-Embodiment：多机器人数据集和 RT-X 模型。投屏帧，约 01:28:56。}
\end{figure}

\begin{knowledgebox}{读图：数据集本身就是基础设施}
图中拼接了大量机器人和任务场景。读这张图要注意，VLA 的瓶颈不只是模型结构，还包括跨实验室、跨硬件、跨任务的数据标准化。
\end{knowledgebox}

\subsection{OpenVLA 与 HiRT：开源和层级动作处理}

OpenVLA 近似开源版 RT-2，把 VLM 接上动作输出，并开源模型与训练流程。HiRT 则指出，直接用 VLM 输出 action 存在动作频率和精细控制问题，因此加入专门处理 Action 的 policy/head：VLM 低频理解和规划，Action Policy 高频执行控制。

\begin{figure}[H]
\centering
\includegraphics[width=0.96\textwidth]{ch25-05531.jpg}
\caption{OpenVLA：开源 VLA 模型，将 VLM 与机器人动作输出连接。投屏帧，约 01:32:11。}
\end{figure}

\begin{figure}[H]
\centering
\includegraphics[width=0.96\textwidth]{ch26-05765.jpg}
\caption{HiRT：通过层级 Robot Transformers 增强机器人控制。投屏帧，约 01:36:05。}
\end{figure}

\begin{knowledgebox}{读图：Action Head 是 VLA 工程化的重要补丁}
HiRT 图中把高层 VLM 信息传给 Action Policy。它对应一个工程事实：大 VLM 推理慢，机器人控制频率高，动作输出需要专门模块来处理连续控制和局部视觉反馈。
\end{knowledgebox}

\subsection{Figure Helix、pi0 与 GR00T N1}

接下来这一段进入企业界和更接近产品化的系统。前面的 OpenVLA 与 HiRT 已经暴露出一个核心张力：VLM 擅长低频理解和语义规划，但真实机器人需要高频、连续、稳定的 action。Figure Helix、pi0 和 GR00T N1 的共同点，是把“理解模块”和“动作模块”更明确地拆开，再用系统设计把两者重新接起来。

Figure Helix 虽未发论文，但代表企业界最新架构：上层系统二像预训练 VLM，下层系统一像实时动作控制。pi0 把 flow model 引入 VLA，用 Action Expert 做通用机器人控制。NVIDIA GR00T N1 则强调人形机器人开放基础模型，也采用 VLM 加 action 处理的结构。

\begin{figure}[H]
\centering
\includegraphics[width=0.92\textwidth]{ch27-05946.jpg}
\caption{Figure AI Helix：系统二与系统一架构示意。投屏帧，约 01:39:06。}
\end{figure}

\begin{figure}[H]
\centering
\includegraphics[width=0.92\textwidth]{ch28-06038.jpg}
\caption{pi0：Vision-Language-Action Flow Model for General Robot Control。投屏帧，约 01:40:38。}
\end{figure}

\begin{figure}[H]
\centering
\includegraphics[width=0.94\textwidth]{ch29-06126.jpg}
\caption{NVIDIA GR00T N1：通用人形机器人的开放基础模型。投屏帧，约 01:42:06。}
\end{figure}

\subsection{本章小结}

RT-2 把 VLM 知识迁移到机器人控制，RT-X/Open X-Embodiment 强调数据规模和跨本体，OpenVLA 提供开源路线，HiRT/pi0/GR00T N1 则开始更认真地处理 Action 模块。这说明 VLA 的核心难点正在从“接上语言和视觉”转向“动作如何高频、稳定、可泛化地输出”。

